\chapter*{ABSTRACT}
\addcontentsline{toc}{chapter}{ABSTRACT}

Reproducibility is a cornerstone of experimental science, yet mobile robotics research
has long struggled to establish it in practice. Setting up a navigation experiment,
recording sensor data across multiple runs, and comparing the resulting trajectories
quantitatively demands a level of infrastructure that most robotics labs build informally,
once per project, with no reuse in mind. This dissertation presents a framework that
makes those steps explicit, automated, and transferable.

The proposed system — called \textit{Fleet UI} — orchestrates autonomous navigation
experiments with one or more mobile robots running ROS~2 Jazzy and the Nav2 navigation
stack. A role-based architecture separates units into \textit{MUUT} (Mobile Unit Under
Test), which executes and records routes, and \textit{FUUT} (Fixed Unit Under Test),
which observes from a fixed position. This separation allows the same physical testbed
to serve fundamentally different experimental designs without changing the software layer.

The record--replay protocol forms the experimental nucleus: a baseline route is defined
by navigating a set of waypoints; the orchestrator stores the visited poses; subsequent
replay runs reproduce those poses using Nav2's \texttt{FollowWaypoints} action.
Sensor bags in MCAP format capture LiDAR, wheel odometry, IMU, and SLAM-estimated pose
at each run. Offline analysis computes pairwise Root Mean Square Error (RMSE) between
trajectories resampled on a normalized time grid, producing trajectory overlays and
repeatability metrics.

Quantitative validation, carried out in Gazebo Harmonic simulation with a single
TurtleBot4 Standard unit, used a clean baseline and ten independent replay executions.
Each replay was started from the same initial state by fully relaunching the simulator,
the Nav2 stack, and the data collection nodes, ensuring that recorded odometry started
at practically the same point in every run. The campaign achieved a mean RMSE of
3.35\,cm between each replay and the reference run, with a maximum of 4.76\,cm, over
trajectories with a mean length of 0.826\,m. The mean final endpoint error was
1.64\,cm (maximum 3.26\,cm), and the mean duration was 17.51\,s with a standard
deviation of 0.36\,s. These results are comfortably below Nav2's default 25\,cm goal
tolerance; the architecture supports multiple units, but empirical evaluation of
simultaneous multi-robot fleet scenarios was not conducted in this dissertation. A web
interface built on React and FastAPI allows researchers to define experiments, monitor
sensor state, and visualise trajectories on the SLAM map without direct exposure to the
ROS~2 command line.

\vspace{0.5cm}
\noindent\textbf{Keywords}: autonomous navigation; path repeatability; mobile robot fleet;
ROS~2; Nav2; SLAM; sensor data collection; experimental framework.
